\documentclass{article}
\usepackage{geometry}
\geometry{
	a4paper,
	left=2.5cm,   
	right=2.5cm,   
	top=2.5cm,      
	bottom=2.5cm,  
	headheight=13.6pt
}

% 最简单、最稳定的中文方案（自动适配系统，永不报错）
\usepackage[UTF8, scheme=plain]{ctex}

% 设置英文默认字体；若系统无 Times New Roman，则自动使用兼容字体
\usepackage{fontspec}
\IfFontExistsTF{Times New Roman}{\setmainfont{Times New Roman}}{\setmainfont{DejaVu Serif}}

\usepackage{amsmath}
\usepackage{amssymb}
\usepackage{booktabs}
\usepackage{caption}
\usepackage{setspace}
\usepackage{array}
\usepackage{listings}
\usepackage{graphicx}
\usepackage{enumitem}
\usepackage{subcaption}
\usepackage{multirow}
\usepackage{xcolor}
\usepackage{colortbl}
\usepackage{nccmath}
\everymath{\displaystyle}  % 全文行内公式自动启用 displaystyle


% 定义代码颜色
\definecolor{codegreen}{rgb}{0,0.6,0}
\definecolor{NavyBlue}{RGB}{0,0,128}
\definecolor{PineGreen}{RGB}{1,121,111}

% 配置 listings 环境
\lstset{
	language=Python,
	tabsize=4,
	captionpos=b,
	numbers=left,
	numbersep=1em,
	sensitive=true,
	showtabs=false,
	frame=shadowbox,
	breaklines=true,
	keepspaces=true,
	showspaces=false,
	showstringspaces=false,
	breakatwhitespace=false,
	basicstyle=\ttfamily\small,
	keywordstyle=\color{NavyBlue},
	commentstyle=\color{codegreen},
	numberstyle=\color{black},
	stringstyle=\color{PineGreen!90!black},
	rulesepcolor=\color{red!20!green!20!blue!20}
}

\definecolor{LightGreen}{rgb}{0.9, 0.95, 0.9}
\definecolor{LightBlue}{RGB}{229, 238, 255}
\definecolor{LightPink}{RGB}{255, 229, 237}

%% 修正表格标题格式：居中对齐，字体为黑体
%\captionsetup[table]{
%	justification=centering,
%	labelsep=quad,
%	font={bf}
%}
%
%% 设置图片标题（图注）字体为黑体
%\captionsetup[figure]{
%	justification=centering,
%	labelsep=quad,
%	font={bf}
%}
%
%% 设置子图标题字体为黑体
%\captionsetup[subfigure]{
%	font={bf}
%}

% 常用记号
\newcommand{\R}{\mathbb{R}}
\newcommand{\Q}{\mathbb{Q}}
\newcommand{\N}{\mathbb{N}}
\newcommand{\calM}{\mathcal{M}}
\newcommand{\eps}{\varepsilon}
\newcommand{\meas}{m}
\newcommand{\mouter}{m^*}
\newcommand{\ind}{\chi}
\newcommand{\aee}{\text{a.e.}}
\newcommand{\diam}{\operatorname{diam}}
\newcommand{\dist}{\operatorname{dist}}
\allowdisplaybreaks
\setstretch{1.2}

\begin{document}
	
	\begin{center}
		{\LARGE \textbf{实变函数研讨课题目汇总}}\\[0.3cm]
%		{\large 姓名 \quad 2024********} \\[0.3cm]\\
		{\large 2026年6月5日}\\[0.3cm]
		{\Large \textbf{实变函数论}}
	\end{center}

\section{第一章：集合与点集}

\subsection{集合运算与集合列极限}

\subsubsection*{题1（《实变函数论》p8 思考题2）}
\paragraph{题目}设
\[
A_1\subset A_2\subset\cdots\subset A_n\subset\cdots,\qquad
B_1\subset B_2\subset\cdots\subset B_n\subset\cdots .
\]
证明
\[
\left(\bigcup_{n=1}^{\infty}A_n\right)\cap
\left(\bigcup_{n=1}^{\infty}B_n\right)
=\bigcup_{n=1}^{\infty}(A_n\cap B_n).
\]

\paragraph{答案}先证 ``$\supset$''。若 $x\in A_n\cap B_n$，则 $x\in\bigcup A_n$ 且 $x\in\bigcup B_n$，故
\[
\bigcup_{n=1}^{\infty}(A_n\cap B_n)
\subset
\left(\bigcup_{n=1}^{\infty}A_n\right)\cap
\left(\bigcup_{n=1}^{\infty}B_n\right).
\]
再证 ``$\subset$''。若
\[
x\in \left(\bigcup_{n=1}^{\infty}A_n\right)\cap
\left(\bigcup_{n=1}^{\infty}B_n\right),
\]
则存在 $i,j\in\N$，使得 $x\in A_i$ 且 $x\in B_j$。令 $N=\max\{i,j\}$，由单调性得 $A_i\subset A_N$，$B_j\subset B_N$，故 $x\in A_N\cap B_N$。因此 $x\in\bigcup_{n=1}^{\infty}(A_n\cap B_n)$。两边相等。

\subsection{映射与基数}

\subsubsection*{题2（《实变函数论》p13 思考题2，按作业修改）}
\paragraph{题目}证明：不存在 $\R$ 上的连续函数 $f$，使得 $f$ 在 $\R\setminus\Q$ 到 $\R\setminus\Q$ 上是一一映射，但在 $\Q$ 上不是一一映射。

\paragraph{答案}反证。设存在这样的连续函数 $f$。因为 $f$ 在 $\Q$ 上不是一一映射，所以存在两个不同的有理数 $q_1<q_2$，使得
\[
f(q_1)=f(q_2).
\]
若 $f$ 在 $[q_1,q_2]$ 上为常数，则任取两个不同无理数 $x,y\in(q_1,q_2)$，有 $f(x)=f(y)$，这与 $f$ 在无理数集上一一矛盾。

若 $f$ 在 $[q_1,q_2]$ 上不为常数，则连续函数在闭区间上取到最大值和最小值。由于两端点函数值相同，必存在一个非退化区间 $J$ 中的每个值都在 $[q_1,q_2]$ 内至少取到两次。又 $\Q\cap[q_1,q_2]$ 可数，故 $f(\Q\cap[q_1,q_2])$ 可数。可在 $J$ 中选取一个数 $c\notin f(\Q\cap[q_1,q_2])$。此时方程 $f(x)=c$ 在 $[q_1,q_2]$ 中至少有两个解，且这些解都不是有理数，于是存在两个不同的无理数 $x,y$ 使 $f(x)=f(y)$，仍与 $f$ 在 $\R\setminus\Q$ 上一一矛盾。因此这样的连续函数不存在。

\subsubsection*{题3（《实变函数论》p17 思考题1）}
\paragraph{题目}设 $A_1\subset A_2$，$B_1\subset B_2$。若 $A_1\sim B_1$，$A_2\sim B_2$，是否一定有
\[
A_2\setminus A_1\sim B_2\setminus B_1?
\]

\paragraph{答案}不一定。取
\[
A_1=\N,\quad A_2=\N\cup\{a\},
\]
其中 $a\notin\N$；再取
\[
B_1=\N,
\qquad B_2=\N\cup\{b,c\},
\]
其中 $b,c\notin\N$ 且 $b\ne c$。显然 $A_1\sim B_1$，且 $A_2,B_2$ 都是可列集，所以 $A_2\sim B_2$。但是
\[
A_2\setminus A_1=\{a\},\qquad B_2\setminus B_1=\{b,c\},
\]
二者基数分别为 $1$ 与 $2$，不对等。因此原命题不成立。

\subsubsection*{题4（《实变函数论》p17 思考题3）}
\paragraph{题目}若 $A\subset B$ 且 $A\sim A\cup C$，证明
\[
B\sim B\cup C.
\]

\paragraph{答案}由 $A\sim A\cup C$，存在从 $A\cup C$ 到 $A$ 的一一映射 $\varphi$。将 $B\cup C$ 分解为
\[
B\cup C=(A\cup C)\cup\bigl(B\setminus(A\cup C)\bigr),
\]
且这两个集合互不相交。定义映射 $F:B\cup C\to B$ 为
\[
F(x)=
\begin{cases}
\varphi(x),&x\in A\cup C,\\
x,&x\in B\setminus(A\cup C).
\end{cases}
\]
因为 $\varphi(A\cup C)=A$，而 $B\setminus(A\cup C)\subset B\setminus A$，两部分像互不相交，所以 $F$ 是从 $B\cup C$ 到
\[
A\cup\bigl(B\setminus(A\cup C)\bigr)\subset B
\]
的一一映射，故 $B\cup C$ 与 $B$ 的某子集对等，即 $|B\cup C|\le |B|$。另一方面 $B\subset B\cup C$，故 $|B|\le |B\cup C|$。由 Cantor--Bernstein 定理，$B\sim B\cup C$。

\subsubsection*{题5（《实变函数论》p25 思考题10）}
\paragraph{题目}设 $E\subset\R^2$ 是不可数集。证明：存在 $x_0\in E$，使得对任一包含 $x_0$ 的圆邻域 $B(x_0,r)$，点集
\[
E\cap B(x_0,r)
\]
是不可数集。

\paragraph{答案}反证。若结论不成立，则对每个 $x\in E$，存在 $r_x>0$，使得 $E\cap B(x,r_x)$ 是可数集。于是 $\{B(x,r_x):x\in E\}$ 是 $E$ 的开覆盖。由 $\R^2$ 中开覆盖的 Lindelöf 性质，可取可数子覆盖
\[
E\subset\bigcup_{j=1}^{\infty}B(x_j,r_{x_j}).
\]
于是
\[
E=\bigcup_{j=1}^{\infty}\bigl(E\cap B(x_j,r_{x_j})\bigr).
\]
右边是可数个可数集的并，因此是可数集，这与 $E$ 不可数矛盾。故存在所求 $x_0$。

\subsubsection*{题6（《实变函数论》p25 思考题11）}
\paragraph{题目}设 $E\subset\R$，且 $|E|<\mathfrak c$，其中 $\mathfrak c=|\R|$ 为连续统的基数。证明：存在实数 $a$，使得
\[
E+\{a\}=\{x+a:x\in E\}\subset \R\setminus\Q .
\]

\paragraph{答案}要使 $E+\{a\}\subset \R\setminus\Q$，等价于对任意 $x\in E$，都有 $x+a\notin\Q$。也就是说，不能存在 $q\in\Q$ 与 $x\in E$ 使得
\[
x+a=q.
\]
这等价于
\[
a\notin \Q-E:=\{q-x:q\in\Q,\,x\in E\}.
\]
下面证明这样的 $a$ 一定存在。

因为 $\Q$ 是可数集，且 $|E|<\mathfrak c$，所以
\[
|\Q\times E|=\max\{\aleph_0,|E|\}<\mathfrak c.
\]
映射
\[
(q,x)\longmapsto q-x
\]
把 $\Q\times E$ 映到 $\Q-E$ 上，因此
\[
|\Q-E|\le |\Q\times E|<\mathfrak c.
\]
而 $|\R|=\mathfrak c$，故 $\Q-E$ 不可能等于整个 $\R$。于是可取
\[
a\in\R\setminus(\Q-E).
\]
对任意 $x\in E$，若 $x+a\in\Q$，则 $a=(x+a)-x\in\Q-E$，与 $a$ 的选取矛盾。因此 $x+a\notin\Q$。故
\[
E+\{a\}\subset \R\setminus\Q.
\]

\subsection{习题1中的点集与基数题}

\subsubsection*{题7（《实变函数论》p54 习题7）}
\paragraph{题目}设 $f(x)$ 是定义在 $[0,1]$ 上的实值函数，且存在常数 $M$，使得对于 $[0,1]$ 中任意有限个数 $x_1,x_2,\ldots,x_n$，均有
\[
\bigl|f(x_1)+f(x_2)+\cdots+f(x_n)\bigr|\le M.
\]
证明集合
\[
E=\{x\in[0,1]:f(x)\ne0\}
\]
是可数集。

\paragraph{答案}对每个正整数 $k$，令
\[
E_k^+=\left\{x\in[0,1]:f(x)\ge \frac1k\right\},\qquad
E_k^- =\left\{x\in[0,1]:f(x)\le -\frac1k\right\}.
\]
先证 $E_k^+$ 是有限集。若 $E_k^+$ 不是有限集，则可在其中取互不相同的 $n$ 个点 $x_1,x_2,\ldots,x_n$，其中 $n>kM$。于是
\[
f(x_1)+f(x_2)+\cdots+f(x_n)\ge \frac{n}{k}>M,
\]
从而
\[
\bigl|f(x_1)+\cdots+f(x_n)\bigr|>M,
\]
与题设矛盾。因此 $E_k^+$ 只能是有限集。

同理，若 $E_k^-$ 不是有限集，则可取互不相同的 $n>kM$ 个点 $y_1,\ldots,y_n\in E_k^-$，于是
\[
f(y_1)+\cdots+f(y_n)\le -\frac{n}{k}<-M,
\]
也与题设矛盾。所以 $E_k^-$ 也是有限集。

最后注意到：若 $f(x)>0$，则存在 $k\in\N$，使得 $f(x)\ge 1/k$；若 $f(x)<0$，则存在 $k\in\N$，使得 $f(x)\le -1/k$。因此
\[
E=\bigcup_{k=1}^{\infty}\left(E_k^+\cup E_k^-\right).
\]
右边是可数个有限集的并，故为可数集。因此 $E$ 是可数集。

\subsubsection*{题6（《实变函数论》p54 习题9）}
\paragraph{题目}设 $E$ 是三维欧氏空间 $\R^3$ 中的点集，且 $E$ 中任意两点的距离都是有理数。证明 $E$ 是可数集。

\paragraph{答案}若 $E$ 有限，结论显然。设 $E$ 无限。从 $E$ 中取一个极大仿射无关点组
\[
a_0,a_1,\ldots,a_d\qquad (0\le d\le 3).
\]
极大性说明 $E$ 全部包含在这些点张成的 $d$ 维仿射子空间中。对每个 $x\in E$，考虑向量
\[
\Phi(x)=\bigl(|x-a_0|,|x-a_1|,\ldots,|x-a_d|\bigr).
\]
由于 $E$ 中任意两点距离均为有理数，故 $\Phi(x)\in\Q^{d+1}$。在 $d$ 维仿射子空间中，一个点由它到 $d+1$ 个仿射无关点的距离唯一确定。因此 $\Phi$ 是单射。于是 $E$ 可嵌入可数集 $\Q^{d+1}$ 中，所以 $E$ 是可数集。

\subsubsection*{题7（《实变函数论》p55 习题11）}
\paragraph{题目}设 $\{f_\alpha(x)\}_{\alpha\in I}$ 是定义在 $[a,b]$ 上的实值函数族。若存在 $M>0$，使得
\[
|f_\alpha(x)|\le M,\qquad x\in[a,b],\ \alpha\in I,
\]
证明：对 $[a,b]$ 中任一可数集 $E$，总存在函数列 $\{f_{\alpha_n}(x)\}$，使得极限
\[
\lim_{n\to\infty}f_{\alpha_n}(x)
\]
在每个 $x\in E$ 上都存在。

\paragraph{答案}将 $E$ 枚举为 $E=\{x_1,x_2,\ldots\}$。任取函数列 $\{f_{\beta_n}\}$，其中 $\beta_n\in I$ 且可视为从函数族中选取。由于 $\{f_{\beta_n}(x_1)\}$ 有界，由 Bolzano--Weierstrass 定理可取子列，使其在 $x_1$ 处收敛；再从该子列中取子列，使其在 $x_2$ 处收敛；依次进行。得到逐次嵌套子列
\[
\{f^{(1)}_n\}\supset\{f^{(2)}_n\}\supset\cdots,
\]
其中第 $k$ 个子列在 $x_1,
\ldots,x_k$ 处均收敛。取对角线列
\[
g_n=f^{(n)}_n.
\]
则对固定的 $x_j$，当 $n\ge j$ 后，$\{g_n(x_j)\}$ 是第 $j$ 个收敛子列的一个尾子列，故收敛。于是 $\{g_n\}$ 即为所求函数列。

\subsubsection*{题8（《实变函数论》p55 习题12）}
\paragraph{题目}设
\[
E=\bigcup_{n=1}^{\infty}A_n,
\]
若 $E$ 的基数为连续统 $\mathfrak c$，证明存在 $n_0$，使得 $A_{n_0}$ 的基数也为 $\mathfrak c$。

\paragraph{答案}若所有 $A_n$ 的基数都小于 $\mathfrak c$，则
\[
|E|\le\sum_{n=1}^{\infty}|A_n|.
\]
连续统 $\mathfrak c=2^{\aleph_0}$ 的共尾数大于 $\aleph_0$，因此可数个基数小于 $\mathfrak c$ 的集合之并仍不能达到基数 $\mathfrak c$。这与 $|E|=\mathfrak c$ 矛盾。故必存在 $n_0$，使得 $|A_{n_0}|=\mathfrak c$。

\subsubsection*{题9（《实变函数论》p55 习题15）}
\paragraph{题目}设 $F\subset\R^n$ 是闭集，$r>0$。证明点集
\[
E=\{t\in\R^n:\text{存在 }x\in F,\ |t-x|=r\}
\]
是闭集。

\paragraph{答案}设 $t_k\in E$ 且 $t_k\to t$。由 $t_k\in E$，存在 $x_k\in F$，使得
\[
|t_k-x_k|=r.
\]
由于 $t_k\to t$，点列 $\{t_k\}$ 有界，于是 $\{x_k\}$ 也有界。取子列 $x_{k_j}\to x$。因为 $F$ 是闭集，故 $x\in F$。又
\[
|t-x|=\lim_{j\to\infty}|t_{k_j}-x_{k_j}|=r.
\]
因此 $t\in E$。所以 $E$ 包含自己的极限点，是闭集。

\subsubsection*{题10（《实变函数论》p55 习题17）}
\paragraph{题目}设 $E\subset\R^2$ 是闭集。对固定的 $y\in\R$，记
\[
E_y=\{x\in\R:(x,y)\in E\}.
\]
证明 $E_y$ 是闭集。

\paragraph{答案}设 $x_k\in E_y$ 且 $x_k\to x$。由 $x_k\in E_y$，有 $(x_k,y)\in E$。又
\[
(x_k,y)\to(x,y).
\]
因为 $E$ 是 $\R^2$ 中的闭集，所以 $(x,y)\in E$，即 $x\in E_y$。因此 $E_y$ 是闭集。

\subsubsection*{题11（《实变函数论》p55 习题19）}
\paragraph{题目}设 $f(x)$ 在 $\R$ 上具有介值性。若对任意 $r\in\Q$，点集
\[
\{x\in\R:f(x)=r\}
\]
都是闭集，证明 $f\in C(\R)$。

\paragraph{答案}反证。若 $f$ 在 $x_0$ 不连续，则存在点列 $x_n\to x_0$ 与 $\eps>0$，使得一列子列满足
\[
f(x_{n_j})\ge f(x_0)+\eps
\]
或
\[
f(x_{n_j})\le f(x_0)-\eps.
\]
只证第一种，第二种类似。取有理数 $r$ 满足
\[
f(x_0)<r<f(x_0)+\eps.
\]
由于 $f$ 有介值性，在 $x_0$ 与 $x_{n_j}$ 之间存在 $y_j$，使得
\[
f(y_j)=r.
\]
显然 $y_j\to x_0$。而集合 $\{x:f(x)=r\}$ 是闭集，所以由 $y_j\in\{f=r\}$ 且 $y_j\to x_0$ 得 $x_0\in\{f=r\}$，即 $f(x_0)=r$，矛盾。故 $f$ 在任意点连续。

\section{第二章：Lebesgue 测度}

\subsection{Lebesgue 外测度与可测集}

\subsubsection*{题12（《实变函数论》p71 思考题3）}
\paragraph{题目}设 $E\subset\R^n$。若对任意 $x\in E$，存在开球 $B(x,r_x)$，使得
\[
\mouter\bigl(E\cap B(x,r_x)\bigr)=0,
\]
证明 $\mouter(E)=0$。

\paragraph{答案}集合族 $\{B(x,r_x):x\in E\}$ 是 $E$ 的一个开覆盖。由 $\R^n$ 的 Lindelöf 性质，可取可数子覆盖
\[
E\subset\bigcup_{j=1}^{\infty}B(x_j,r_{x_j}).
\]
于是
\[
E=\bigcup_{j=1}^{\infty}\bigl(E\cap B(x_j,r_{x_j})\bigr).
\]
由外测度的次可加性，
\[
\mouter(E)
\le\sum_{j=1}^{\infty}\mouter\bigl(E\cap B(x_j,r_{x_j})\bigr)=0.
\]
故 $\mouter(E)=0$。

\subsubsection*{题13（《实变函数论》p78 思考题2）}
\paragraph{题目}设 $\{A_n\}$ 是两两互不相交的可测集列，且 $B_n\subset A_n$，$n=1,2,\ldots$。证明
\[
\mouter\left(\bigcup_{n=1}^{\infty}B_n\right)=\sum_{n=1}^{\infty}\mouter(B_n).
\]

\paragraph{答案}先对有限并证明。因 $A_1,\ldots,A_N$ 可测且互不相交，利用可测集对任意集合的分割可得
\[
\mouter\left(\bigcup_{n=1}^{N}B_n\right)=\sum_{n=1}^{N}\mouter(B_n).
\]
于是由单调性，
\[
\mouter\left(\bigcup_{n=1}^{\infty}B_n\right)
\ge
\mouter\left(\bigcup_{n=1}^{N}B_n\right)
=
\sum_{n=1}^{N}\mouter(B_n).
\]
令 $N\to\infty$，得
\[
\mouter\left(\bigcup_{n=1}^{\infty}B_n\right)
\ge
\sum_{n=1}^{\infty}\mouter(B_n).
\]
反向不等式由外测度的次可加性直接得到：
\[
\mouter\left(\bigcup_{n=1}^{\infty}B_n\right)
\le
\sum_{n=1}^{\infty}\mouter(B_n).
\]
两式合并即得结论。

\subsubsection*{题14（《实变函数论》p78 思考题3）}
\paragraph{题目}设 $E_1,E_2\subset\R^n$，且 $E_1$ 是可测集。若
\[
\meas(E_1\triangle E_2)=0,
\]
证明 $E_2$ 是可测集，且
\[
\meas(E_2)=\meas(E_1).
\]

\paragraph{答案}因为 $E_1\triangle E_2$ 是零测集，所以它的任意子集都是可测零测集。注意
\[
E_2=E_1\triangle(E_1\triangle E_2).
\]
右边是两个可测集的对称差，因此 $E_2$ 可测。又
\[
E_2\setminus E_1\subset E_1\triangle E_2,
\qquad
E_1\setminus E_2\subset E_1\triangle E_2,
\]
故
\[
\meas(E_2\setminus E_1)=\meas(E_1\setminus E_2)=0.
\]
于是
\[
\meas(E_2)=\meas(E_1\cap E_2)+\meas(E_2\setminus E_1)
=\meas(E_1\cap E_2)
=\meas(E_1).
\]

\subsubsection*{题15（《实变函数论》p78 思考题4）}
\paragraph{题目}设点集 $B$ 满足：对任给 $\eps>0$，都存在可测集 $A$，使得
\[
\mouter(A\triangle B)<\eps.
\]
证明 $B$ 是可测集。

\paragraph{答案}任取 $T\subset\R^n$。由外测度的次可加性，恒有
\[
\mouter(T)\le \mouter(T\cap B)+\mouter(T\cap B^c).
\]
反向估计如下。取可测集 $A$ 使 $\mouter(A\triangle B)<\eps$。因为 $A$ 可测，
\[
\mouter(T)=\mouter(T\cap A)+\mouter(T\cap A^c).
\]
又
\[
T\cap A\subset (T\cap B)\cup(A\triangle B),
\]
以及
\[
T\cap A^c\subset (T\cap B^c)\cup(A\triangle B),
\]
故
\[
\mouter(T)
\le \mouter(T\cap B)+\mouter(T\cap B^c)+2\mouter(A\triangle B)
<\mouter(T\cap B)+\mouter(T\cap B^c)+2\eps.
\]
令 $\eps\to0$，得到
\[
\mouter(T)\le \mouter(T\cap B)+\mouter(T\cap B^c).
\]
因此
\[
\mouter(T)= \mouter(T\cap B)+\mouter(T\cap B^c),
\]
由 Carath\'eodory 条件，$B$ 可测。

\subsubsection*{题16（《实变函数论》p84 思考题1）}
\paragraph{题目}设 $E\subset\R^n$，且 $\mouter(E)<+\infty$。若
\[
\mouter(E)=\sup\{\meas(F):F\subset E,\ F\text{ 是有界闭集}\},
\]
证明 $E$ 是可测集。

\paragraph{答案}对每个 $k\in\N$，取有界闭集 $F_k\subset E$，使得
\[
\meas(F_k)>\mouter(E)-\frac1k.
\]
因 $F_k$ 是闭集，故可测。由 Carath\'eodory 条件作用于可测集 $F_k$ 和任意集合 $E$，有
\[
\mouter(E)=\meas(F_k)+\mouter(E\setminus F_k).
\]
于是
\[
\mouter(E\setminus F_k)<\frac1k.
\]
令
\[
F=\bigcup_{k=1}^{\infty}F_k.
\]
则 $F$ 可测且 $F\subset E$。又对每个 $k$，$E\setminus F\subset E\setminus F_k$，所以
\[
\mouter(E\setminus F)\le\mouter(E\setminus F_k)<\frac1k.
\]
令 $k\to\infty$，得 $\mouter(E\setminus F)=0$。因此
\[
E=F\cup(E\setminus F)
\]
是可测集与零测集的并，故 $E$ 可测。

\subsubsection*{题17（《实变函数论》p84 思考题2）}
\paragraph{题目}设 $A$ 是可测集，$B\subset\R^n$。证明
\[
\mouter(A\cup B)+\mouter(A\cap B)=\meas(A)+\mouter(B).
\]

\paragraph{答案}因为 $A$ 可测，对集合 $A\cup B$ 作分割，得
\[
\mouter(A\cup B)
=\mouter((A\cup B)\cap A)+\mouter((A\cup B)\cap A^c)
=\meas(A)+\mouter(B\setminus A).
\]
又对集合 $B$ 作分割，得
\[
\mouter(B)=\mouter(B\cap A)+\mouter(B\cap A^c)
=\mouter(A\cap B)+\mouter(B\setminus A).
\]
两式相减或代入，即得
\[
\mouter(A\cup B)+\mouter(A\cap B)=\meas(A)+\mouter(B).
\]

\subsection{习题2中的测度题}

\subsubsection*{题18（《实变函数论》p94 习题8）}
\paragraph{题目}设 $\{E_k\}$ 是 $[0,1]$ 中的可测集列，$\meas(E_k)=1$，$k=1,2,\ldots$。证明
\[
\meas\left(\bigcap_{k=1}^{\infty}E_k\right)=1.
\]

\paragraph{答案}因为 $E_k\subset[0,1]$ 且 $\meas(E_k)=1$，所以
\[
\meas([0,1]\setminus E_k)=0.
\]
于是
\[
[0,1]\setminus\bigcap_{k=1}^{\infty}E_k
=\bigcup_{k=1}^{\infty}([0,1]\setminus E_k).
\]
由可数次可加性，右边为零测集。因此
\[
\meas\left([0,1]\setminus\bigcap_{k=1}^{\infty}E_k\right)=0,
\]
从而
\[
\meas\left(\bigcap_{k=1}^{\infty}E_k\right)=1.
\]

\subsubsection*{题19（《实变函数论》p94 习题9）}
\paragraph{题目}设 $E_1,E_2,\ldots,E_k$ 是 $[0,1]$ 中的可测集，且
\[
\sum_{i=1}^{k}\meas(E_i)>k-1.
\]
证明
\[
\meas\left(\bigcap_{i=1}^{k}E_i\right)>0.
\]

\paragraph{答案}由 De Morgan 公式，
\[
[0,1]\setminus\bigcap_{i=1}^{k}E_i
=\bigcup_{i=1}^{k}([0,1]\setminus E_i).
\]
因此
\[
\meas\left([0,1]\setminus\bigcap_{i=1}^{k}E_i\right)
\le \sum_{i=1}^{k}\meas([0,1]\setminus E_i)
=\sum_{i=1}^{k}(1-\meas(E_i))
=k-\sum_{i=1}^{k}\meas(E_i)<1.
\]
所以
\[
\meas\left(\bigcap_{i=1}^{k}E_i\right)>0.
\]

\subsubsection*{题20（《实变函数论》p95 习题11）}
\paragraph{题目}设 $\{B_\alpha\}_{\alpha\in I}$ 是 $\R^n$ 中的一族开球，记
\[
G=\bigcup_{\alpha\in I}B_\alpha.
\]
若 $0<\lambda<\meas(G)$，证明存在有限个互不相交的开球
\[
B_{\alpha_1},B_{\alpha_2},\ldots,B_{\alpha_N}
\]
使得
\[
\sum_{j=1}^{N}\meas(B_{\alpha_j})>\frac{\lambda}{3^n}.
\]

\paragraph{答案}取紧集 $K\subset G$，使得 $\meas(K)>\lambda$。由于 $\{B_\alpha\}$ 覆盖 $K$，由 Heine--Borel 定理可取有限子覆盖
\[
K\subset\bigcup_{j=1}^{M}B_j.
\]
在这有限个球中用贪心法选球：先取直径最大的一个球 $B_{\alpha_1}$，删去所有与它相交的球；在剩下的球中再取直径最大的球 $B_{\alpha_2}$，继续下去，直到没有球剩余。所得 $B_{\alpha_1},\ldots,B_{\alpha_N}$ 两两不相交。

任一未被选中的球都与某个已选球相交，且其半径不大于该已选球的半径，因此它包含在对应已选球同心半径扩大三倍的球 $3B_{\alpha_j}$ 中。于是
\[
K\subset\bigcup_{j=1}^{N}3B_{\alpha_j}.
\]
所以
\[
\lambda<\meas(K)
\le \sum_{j=1}^{N}\meas(3B_{\alpha_j})
=3^n\sum_{j=1}^{N}\meas(B_{\alpha_j}).
\]
即
\[
\sum_{j=1}^{N}\meas(B_{\alpha_j})>\frac{\lambda}{3^n}.
\]

\subsubsection*{题21（《实变函数论》p95 习题12）}
\paragraph{题目}设 $\{B_k\}$ 是 $\R^n$ 中递减可测集列，$\mouter(A)<+\infty$。令
\[
E_k=A\cap B_k,
\qquad
E=\bigcap_{k=1}^{\infty}E_k.
\]
证明
\[
\lim_{k\to\infty}\mouter(E_k)=\mouter(E).
\]

\paragraph{答案}记 $C_k=B_1\setminus B_k$，则 $\{C_k\}$ 是递增可测集列，并且
\[
C=\bigcup_{k=1}^{\infty}C_k=B_1\setminus\bigcap_{k=1}^{\infty}B_k.
\]
由于 $B_k$ 可测，$C_k$ 可测。对集合 $A\cap B_1$ 用可测集 $C_k$ 分割，得到
\[
\mouter(A\cap B_1)=\mouter(A\cap C_k)+\mouter(A\cap B_k).
\]
同理，
\[
\mouter(A\cap B_1)=\mouter(A\cap C)+\mouter\left(A\cap\bigcap_{k=1}^{\infty}B_k\right).
\]
又 $A\cap C_k$ 递增到 $A\cap C$，且外测度有限，故
\[
\lim_{k\to\infty}\mouter(A\cap C_k)=\mouter(A\cap C).
\]
两式相减即得
\[
\lim_{k\to\infty}\mouter(A\cap B_k)
=
\mouter\left(A\cap\bigcap_{k=1}^{\infty}B_k\right),
\]
即
\[
\lim_{k\to\infty}\mouter(E_k)=\mouter(E).
\]

\subsubsection*{题22（《实变函数论》p95 习题13）}
\paragraph{题目}设 $E\subset\R^n$，$H\supset E$ 且 $H$ 是可测集。若 $H\setminus E$ 的任一可测子集皆为零测集，试问：$H$ 是 $E$ 的等测包吗？

\paragraph{答案}是。若 $H$ 不是 $E$ 的等测包，则有
\[
\meas(H)>\mouter(E).
\]
取 $\eps>0$ 使得
\[
\mouter(E)+\eps<\meas(H).
\]
由外测度定义，可取开集 $G\supset E$，使得
\[
\meas(G)<\mouter(E)+\eps<\meas(H).
\]
于是 $H\setminus G$ 是可测集，且
\[
H\setminus G\subset H\setminus E.
\]
另一方面，
\[
\meas(H\setminus G)
\ge \meas(H)-\meas(G)>0.
\]
这就得到 $H\setminus E$ 中存在正测度可测子集，矛盾。因此 $\meas(H)=\mouter(E)$，即 $H$ 是 $E$ 的等测包。

\subsubsection*{题23（《实变函数论》p95 习题15）}
\paragraph{题目}设 $E\subset[0,1]$ 是可测集，且 $\meas(E)\ge c>0$。给定 $x_i\in[0,1]$，$i=1,2,\ldots,n$，其中
\[
n>\frac{2}{c}.
\]
证明：$E$ 中存在两个点，其距离等于 $\{x_1,x_2,\ldots,x_n\}$ 中某两个点之间的距离。

\paragraph{答案}反证。假设 $E$ 中任意两点的距离都不等于 $|x_i-x_j|$，$i\ne j$。考虑平移集
\[
E+x_i=\{y+x_i:y\in E\},\qquad i=1,2,\ldots,n.
\]
若 $(E+x_i)\cap(E+x_j)\ne\varnothing$，则存在 $u,v\in E$，使得
\[
u+x_i=v+x_j,
\]
从而
\[
|u-v|=|x_i-x_j|,
\]
与假设矛盾。因此 $E+x_1,\ldots,E+x_n$ 两两不交。它们都包含于 $[0,2]$，于是
\[
2\ge \meas\left(\bigcup_{i=1}^{n}(E+x_i)\right)
=\sum_{i=1}^{n}\meas(E+x_i)=n\meas(E)
\ge nc>2,
\]
矛盾。故结论成立。

\section{第三章：可测函数}

\subsection{可测函数的定义与连续函数关系}

\subsubsection*{题24（3.12第一次作业）}
\paragraph{题目}设 $f:E\to\R$，记
\[
E[f>\alpha]=\{x\in E:f(x)>\alpha\},
\qquad
E[f\ge\alpha]=\{x\in E:f(x)\ge\alpha\}.
\]
证明
\[
E[f\ge\alpha]=\bigcap_{n=1}^{\infty}E\left[f>\alpha-\frac1n\right]
=\bigcap_{n=1}^{\infty}E\left[f\ge\alpha-\frac1n\right].
\]

\paragraph{答案}若 $x\in E[f\ge\alpha]$，则 $f(x)\ge\alpha>\alpha-1/n$，故 $x\in E[f>\alpha-1/n]$，也属于 $E[f\ge\alpha-1/n]$，对一切 $n$ 成立。

反过来，若 $x\in\bigcap_{n=1}^{\infty}E[f>\alpha-1/n]$，则对任意 $n$，
\[
f(x)>\alpha-\frac1n.
\]
令 $n\to\infty$ 得 $f(x)\ge\alpha$。因此
\[
\bigcap_{n=1}^{\infty}E\left[f>\alpha-\frac1n\right]
\subset E[f\ge\alpha].
\]
同理，若 $f(x)\ge\alpha-1/n$ 对一切 $n$ 成立，令 $n\to\infty$ 也得 $f(x)\ge\alpha$。故两个交集都等于 $E[f\ge\alpha]$。

\subsubsection*{题25（3.12第一次作业）}
\paragraph{题目}设 $\R^n$ 上的实值函数列 $\{f_k\}$ 满足 $f_k(x)\to f(x)$。证明对任意 $\alpha\in\R$，
\[
\{x\in\R^n:f(x)>\alpha\}
=
\bigcup_{j=1}^{\infty}\bigcup_{m=1}^{\infty}
\bigcap_{k=m}^{\infty}
\left\{x\in\R^n:f_k(x)\ge \alpha+\frac1j\right\}.
\]

\paragraph{答案}若 $f(x)>\alpha$，可取 $j$ 足够大，使得
\[
f(x)>\alpha+\frac1j.
\]
由 $f_k(x)\to f(x)$，存在 $m$，使得当 $k\ge m$ 时，
\[
f_k(x)\ge\alpha+\frac1j.
\]
故 $x$ 属于右边。

反过来，若 $x$ 属于右边，则存在 $j,m$，使得对所有 $k\ge m$，
\[
f_k(x)\ge\alpha+\frac1j.
\]
令 $k\to\infty$，得
\[
f(x)\ge\alpha+\frac1j>\alpha.
\]
故 $x\in\{f>\alpha\}$。两边相等。

\subsubsection*{题26（3.19第二次作业）}
\paragraph{题目}设 $f(x)$ 是定义在 $\R$ 上的实值函数。证明点集
\[
E=\left\{x\in\R:\lim_{y\to x}f(y)=+\infty\right\}
\]
是可数集。

\paragraph{答案}对每个 $m\in\N$，令
\[
E_m=\{x\in E: |f(x)|\le m\}.
\]
则
\[
E=\bigcup_{m=1}^{\infty}E_m.
\]
只需证明每个 $E_m$ 至多可数。

固定 $m$。若 $x\in E_m$，由 $\lim_{y\to x}f(y)=+\infty$，存在 $\delta_x>0$，使得当 $0<|y-x|<\delta_x$ 时，
\[
f(y)>m+1.
\]
若存在 $z\in E_m$，$z\ne x$，且 $|z-x|<\delta_x$，则由 $z\in E_m$ 得 $|f(z)|\le m$，特别 $f(z)\le m$；但由 $0<|z-x|<\delta_x$ 得 $f(z)>m+1$，矛盾。因此每个 $x\in E_m$ 都是 $E_m$ 的孤立点。

实直线中的孤立点集至多可数：给每个孤立点选一个含它而不含其他点的有理端点开区间，便得到到可数有理区间族的单射。因此 $E_m$ 至多可数，从而 $E$ 是可数个可数集的并，故 $E$ 可数。

\subsubsection*{题27（3.19第二次作业）}
\paragraph{题目}设 $(a,b)$ 中全体有理数为 $\{r_n\}_{n\in\N}$，并且
\[
\sum_{n=1}^{\infty}C_n<+\infty,
\qquad C_n>0.
\]
定义 $(a,b)$ 上的函数
\[
f(x)=\sum_{r_n<x}C_n.
\]
证明：$f$ 是递增函数；$f$ 在 $(a,b)$ 上的无理点连续，在有理点间断；并且
\[
f(r_n+)-f(r_n-)=C_n,
\qquad n=1,2,\ldots.
\]

\paragraph{答案}若 $x_1<x_2$，则
\[
\{n:r_n<x_1\}\subset\{n:r_n<x_2\},
\]
且 $C_n>0$，故 $f(x_1)\le f(x_2)$，即 $f$ 递增。

先证无理点连续。设 $x_0$ 是无理数。给定 $\eps>0$，取 $N$ 使得
\[
\sum_{n>N}C_n<\eps.
\]
有限集 $\{r_1,
\ldots,r_N\}$ 中没有点等于 $x_0$，故可取 $\delta>0$，使得区间 $(x_0-\delta,x_0+\delta)$ 不含 $r_1,
\ldots,r_N$。若 $|x-x_0|<\delta$，则 $f(x)-f(x_0)$ 的变化只可能来自指标 $n>N$ 的项，因此
\[
|f(x)-f(x_0)|\le\sum_{n>N}C_n<\eps.
\]
故 $f$ 在 $x_0$ 连续。

再看有理点 $r_m$。由定义可得
\[
f(r_m-)=\sum_{r_n<r_m}C_n,
\qquad
f(r_m+)=\sum_{r_n\le r_m}C_n.
\]
因为枚举中 $r_m$ 只出现一次，所以
\[
f(r_m+)-f(r_m-)=C_m>0.
\]
因此 $f$ 在每个有理点 $r_m$ 处有跳跃间断。

\subsubsection*{题28（3.19第二次作业：Cantor--Bernstein 定理的另一种证明）}
\paragraph{题目}设 $A,B$ 为两个非空集合。若存在 $A_1\subset A$、$B_1\subset B$，以及一一映射
\[
\varphi_1:A\to B_1,
\qquad
\varphi_2:B\to A_1.
\]
记
\[
\varphi=\varphi_2\circ\varphi_1,
\]
并令
\[
A_0=A,
\qquad
A_{2n+2}=\varphi(A_{2n}),
\qquad
A_{2n+3}=\varphi(A_{2n+1}),
\quad n=0,1,2,\ldots.
\]
证明：
\begin{enumerate}[label=(\arabic*)]
\item
\[
A=\left[\bigcup_{n=0}^{\infty}(A_n\setminus A_{n+1})\right]
\cup \bigcap_{n=0}^{\infty}A_n,
\]
且
\[
A_1=\left[\bigcup_{n=1}^{\infty}(A_n\setminus A_{n+1})\right]
\cup \bigcap_{n=1}^{\infty}A_n.
\]
\item 对任意 $n=0,1,2,
\ldots$，有
\[
A_n\setminus A_{n+1}\sim A_{n+2}\setminus A_{n+3},
\]
并由此证明 $A\sim B$。
\end{enumerate}

\paragraph{答案}由于 $\varphi:A\to A_1$ 是一一映射，且 $A_1\subset A$，可知
\[
A=A_0\supset A_1\supset A_2\supset\cdots.
\]
对任意递减集合列都有分解
\[
A_0=\left[\bigcup_{n=0}^{\infty}(A_n\setminus A_{n+1})\right]\cup\bigcap_{n=0}^{\infty}A_n.
\]
同理，从 $A_1\supset A_2\supset\cdots$ 出发，得
\[
A_1=\left[\bigcup_{n=1}^{\infty}(A_n\setminus A_{n+1})\right]\cup\bigcap_{n=1}^{\infty}A_n.
\]
这证明了 (1)。

对 (2)，因为 $\varphi$ 是一一映射，且 $A_{n+2}=\varphi(A_n)$、$A_{n+3}=\varphi(A_{n+1})$，所以
\[
\varphi(A_n\setminus A_{n+1})=A_{n+2}\setminus A_{n+3}.
\]
故
\[
A_n\setminus A_{n+1}\sim A_{n+2}\setminus A_{n+3}.
\]

现在比较 $A$ 与 $A_1$ 的分解：
\[
A=(A_0\setminus A_1)\cup(A_1\setminus A_2)\cup(A_2\setminus A_3)\cup\cdots\cup\bigcap_{n=0}^{\infty}A_n,
\]
而
\[
A_1=(A_1\setminus A_2)\cup(A_2\setminus A_3)\cup\cdots\cup\bigcap_{n=1}^{\infty}A_n.
\]
利用上面得到的等势关系，可以把 $A_0\setminus A_1$ 与 $A_2\setminus A_3$ 配对，把 $A_1\setminus A_2$ 与 $A_3\setminus A_4$ 配对，依次后移；而
\[
\bigcap_{n=0}^{\infty}A_n=\bigcap_{n=1}^{\infty}A_n.
\]
因此 $A\sim A_1$。又 $\varphi_2:B\to A_1$ 是一一映射到 $A_1$ 上，所以 $B\sim A_1$。由传递性得
\[
A\sim B.
\]

\subsubsection*{题29（《实变函数论》p109 思考题7）}
\paragraph{题目}设 $f(x)$ 是 $\R$ 上几乎处处连续的函数。问是否一定存在 $g\in C(\R)$，使得
\[
g(x)=f(x),\qquad \aee\ x\in\R?
\]

\paragraph{答案}不一定。取
\[
f(x)=\ind_{(0,+\infty)}(x)=
\begin{cases}
0,&x\le 0,\\
1,&x>0.
\end{cases}
\]
该函数只在 $x=0$ 处不连续，因此几乎处处连续。若存在连续函数 $g$ 使 $g=f$ a.e.，则在 $(-\infty,0)$ 中，由于 $g$ 连续且 $g=0$ a.e.，可推出 $g\equiv0$；在 $(0,+\infty)$ 中同理推出 $g\equiv1$。于是
\[
\lim_{x\to0^-}g(x)=0,
\qquad
\lim_{x\to0^+}g(x)=1,
\]
与 $g$ 在 $0$ 处连续矛盾。因此这样的 $g$ 不一定存在。

\subsubsection*{题30（《实变函数论》p119 思考题6）}
\paragraph{题目}设 $E\subset\R$ 上的可测函数列 $\{f_k(x)\}$ 满足
\[
f_k(x)\le f_{k+1}(x),\qquad k=1,2,\ldots.
\]
若 $f_k$ 在 $E$ 上依测度收敛到 $0$，问 $f_k$ 是否几乎处处收敛到 $0$？

\paragraph{答案}是。由依测度收敛可知，存在子列 $f_{k_j}$，使得
\[
f_{k_j}(x)\to0,
\qquad \aee\ x\in E.
\]
在这个满测度集合上，固定 $x$。由于原序列关于 $k$ 单调递增，对于 $k_j\le k\le k_{j+1}$，有
\[
f_{k_j}(x)\le f_k(x)\le f_{k_{j+1}}(x).
\]
令 $j\to\infty$，左右两端都趋于 $0$，由夹逼定理得 $f_k(x)\to0$。所以 $f_k\to0$ a.e. 于 $E$。

\subsection{习题3中的可测函数列}

\subsubsection*{题31（《实变函数论》p126 习题2）}
\paragraph{题目}设 $z=f(x,y)$ 是 $\R^2$ 上的连续函数，$g_1(x),g_2(x)$ 是 $[a,b]$ 上的实值可测函数。证明
\[
F(x)=f(g_1(x),g_2(x))
\]
是 $[a,b]$ 上的可测函数。

\paragraph{答案}令
\[
G(x)=(g_1(x),g_2(x)).
\]
若 $U\subset\R^2$ 是开集，则 $U$ 可表示为可数个开矩形的并，而
\[
G^{-1}(I_1\times I_2)=g_1^{-1}(I_1)\cap g_2^{-1}(I_2)
\]
是可测集。因此 $G$ 是从 $[a,b]$ 到 $\R^2$ 的可测映射。对任意开集 $V\subset\R$，由 $f$ 连续知 $f^{-1}(V)$ 是 $\R^2$ 中开集，从而
\[
F^{-1}(V)=G^{-1}(f^{-1}(V))
\]
是可测集。故 $F$ 可测。

\subsubsection*{题32（《实变函数论》p126 习题6）}
\paragraph{题目}设 $\{f_k(x)\}$ 是 $E\subset\R^n$ 上的实值可测函数列，$\meas(E)<+\infty$。证明 $f_k(x)\to0$ a.e. 的充分必要条件是：对任意 $d>0$，有
\[
\lim_{l\to\infty}\meas\left(\left\{x\in E:\sup_{k\ge l}|f_k(x)|\ge d\right\}\right)=0.
\]

\paragraph{答案}记
\[
A_l(d)=\left\{x\in E:\sup_{k\ge l}|f_k(x)|\ge d\right\}.
\]
则 $A_l(d)$ 随 $l$ 递减，并且
\[
\bigcap_{l=1}^{\infty}A_l(d)=\left\{x\in E:\limsup_{k\to\infty}|f_k(x)|\ge d\right\}.
\]
由于 $\meas(E)<+\infty$，由测度的上连续性，
\[
\lim_{l\to\infty}\meas(A_l(d))
=
\meas\left(\bigcap_{l=1}^{\infty}A_l(d)\right).
\]
若 $f_k\to0$ a.e.，则对每个 $d>0$，右边集合是零测集，故极限为 $0$。

反过来，若对每个 $d>0$ 都有上述极限为 $0$，则对每个正整数 $m$，
\[
\meas\left(\left\{x:\limsup_{k\to\infty}|f_k(x)|\ge \frac1m\right\}\right)=0.
\]
取可数并得
\[
\meas\left(\left\{x:\limsup_{k\to\infty}|f_k(x)|>0\right\}\right)=0.
\]
故 $\limsup |f_k(x)|=0$ a.e.，即 $f_k(x)\to0$ a.e.。

\subsubsection*{题33（《实变函数论》p127 习题10）}
\paragraph{题目}设 $f_n(x)$，$n=1,2,\ldots$，是 $[0,1]$ 上的递增函数，并且 $\{f_n(x)\}$ 在 $[0,1]$ 上依测度收敛于 $f(x)$。证明在 $f(x)$ 的连续点 $x_0$ 处，有
\[
f_n(x_0)\to f(x_0).
\]

\paragraph{答案}反证。设存在 $\eps>0$ 和子列 $f_{n_j}$，使得
\[
f_{n_j}(x_0)\ge f(x_0)+\eps
\]
对无穷多个 $j$ 成立。由于 $f$ 在 $x_0$ 连续，可取 $\delta>0$，使得当 $x\in[x_0,x_0+\delta]\cap[0,1]$ 时，
\[
f(x)<f(x_0)+\frac{\eps}{2}.
\]
又每个 $f_{n_j}$ 关于 $x$ 递增，所以在该区间上
\[
f_{n_j}(x)\ge f_{n_j}(x_0)
\ge f(x_0)+\eps>f(x)+\frac{\eps}{2}.
\]
于是
\[
\meas\left\{x:|f_{n_j}(x)-f(x)|>\frac{\eps}{2}\right\}\ge\delta
\]
对无穷多个 $j$ 成立，矛盾于依测度收敛。若有子列满足 $f_{n_j}(x_0)\le f(x_0)-\eps$，则在 $[x_0-\delta,x_0]\cap[0,1]$ 上同理得到矛盾。因此 $f_n(x_0)\to f(x_0)$。

\subsubsection*{题34（《实变函数论》p127 习题14）}
\paragraph{题目}设有定义在可测集 $E\subset\R^n$ 上的函数 $f(x)$，且对任给的 $\alpha>0$，存在 $E$ 中的闭集 $F$，$\meas(E\setminus F)<\alpha$，使得 $f(x)$ 在 $F$ 上连续。证明 $f(x)$ 是 $E$ 上的可测函数。

\paragraph{答案}对每个 $k\in\N$，取闭集 $F_k\subset E$，使得
\[
\meas(E\setminus F_k)<\frac1k,
\]
且 $f$ 在 $F_k$ 上连续。令
\[
F=\bigcup_{k=1}^{\infty}F_k,
\qquad
N=E\setminus F.
\]
由于 $N\subset E\setminus F_k$ 对所有 $k$ 成立，故
\[
\meas(N)=0.
\]
对任意开集 $G\subset\R$，有
\[
f^{-1}(G)=\left[\bigcup_{k=1}^{\infty}(F_k\cap f^{-1}(G))\right]\cup(N\cap f^{-1}(G)).
\]
因为 $f|_{F_k}$ 连续，$F_k\cap f^{-1}(G)$ 是 $F_k$ 中相对开集，故为 Borel 集；而 $N\cap f^{-1}(G)$ 是零测集的子集，故可测。于是 $f^{-1}(G)$ 可测，所以 $f$ 可测。

\subsubsection*{题35（《实变函数论》p127 习题15）}
\paragraph{题目}设 $\{f_n(x)\}$ 是 $[a,b]$ 上的可测函数列，$f(x)$ 是 $[a,b]$ 上的实值函数。若对任给 $d>0$，都有
\[
\lim_{n\to\infty}\mouter\left(\{x\in[a,b]:|f_n(x)-f(x)|>d\}\right)=0,
\]
问 $f(x)$ 是否是 $[a,b]$ 上的可测函数？

\paragraph{答案}是。由条件可取子列 $f_{n_j}$，使得
\[
\mouter\left(\left\{x:|f_{n_j}(x)-f(x)|>\frac1j\right\}\right)<2^{-j}.
\]
对每个 $j$，取可测集 $G_j$ 覆盖
\[
\left\{x:|f_{n_j}(x)-f(x)|>\frac1j\right\}
\]
且 $\meas(G_j)<2^{-j+1}$。令
\[
H_m=\bigcup_{j\ge m}G_j,
\qquad
H=\bigcap_{m=1}^{\infty}H_m.
\]
则
\[
\meas(H)\le \lim_{m\to\infty}\sum_{j\ge m}2^{-j+1}=0.
\]
若 $x\notin H$，则存在 $m$，使得 $x\notin G_j$ 对所有 $j\ge m$ 成立，从而
\[
|f_{n_j}(x)-f(x)|\le\frac1j,
\qquad j\ge m.
\]
所以 $f_{n_j}(x)\to f(x)$ 在 $[a,b]\setminus H$ 上逐点成立。于是 $f$ 在零测集外是可测函数列的逐点极限，故 $f$ 可测。

\subsubsection*{题36（《实变函数论》p127 习题16）}
\paragraph{题目}设 $f(x),f_k(x)$，$k=1,2,\ldots$，是 $E\subset\R$ 上实值可测函数。若对任给 $d>0$，必有
\[
\lim_{n\to\infty}\meas\left(\bigcup_{k=n}^{\infty}\{x:|f_k(x)-f(x)|>d\}\right)=0,
\]
证明对任给 $\eta>0$，存在 $e\subset E$ 且 $\meas(e)<\eta$，使得 $f_k(x)$ 在 $E\setminus e$ 上一致收敛于 $f(x)$。

\paragraph{答案}对每个 $m\in\N$，由条件可取 $N_m$，使得
\[
\meas\left(\bigcup_{k=N_m}^{\infty}\left\{x:|f_k(x)-f(x)|>\frac1m\right\}\right)<\frac{\eta}{2^m}.
\]
令
\[
e=\bigcup_{m=1}^{\infty}\bigcup_{k=N_m}^{\infty}\left\{x:|f_k(x)-f(x)|>\frac1m\right\}.
\]
则
\[
\meas(e)\le\sum_{m=1}^{\infty}\frac{\eta}{2^m}=\eta.
\]
若 $x\in E\setminus e$，则对每个 $m$，当 $k\ge N_m$ 时，
\[
|f_k(x)-f(x)|\le\frac1m.
\]
这正说明 $f_k\to f$ 在 $E\setminus e$ 上一致收敛。

\subsubsection*{题37（5.27作业P118例3）}
\paragraph{题目}设 $f(x),f_k(x)$，$k\in\N$，是 $E\subset\R$ 上的实值可测函数，$\meas(E)<+\infty$。
\begin{enumerate}[label=(\roman*)]
\item 若在任一子列 $\{f_{k_i}(x)\}$ 中均有子列 $\{f_{k_{i_j}}(x)\}$ 在 $E$ 上收敛于 $f(x)$，则 $f_k(x)$ 在 $E$ 上依测度收敛于 $f(x)$。
\item 若 $f_k(x)>0$，且 $f_k(x)$ 在 $E$ 上依测度收敛于 $f(x)$，则对 $p>0$，$f_k^p(x)$ 在 $E$ 上依测度收敛于 $f^p(x)$。
\end{enumerate}

\paragraph{答案}
(i) 反证法。若结论不真，则存在 $\eps_0>0$、$\delta_0>0$ 以及子列 $\{f_{k_i}\}$，使得
\[
\meas\{x\in E:|f_{k_i}(x)-f(x)|>\eps_0\}\ge\delta_0.
\]
但依题设知，在 $\{f_{k_i}\}$ 中存在子列 $\{f_{k_{i_j}}\}$ 使得 $f_{k_{i_j}}(x)$ 在 $E$ 上依测度收敛于 $f(x)$，这与上式矛盾。故 $f_k\to f$ 依测度成立。

(ii) 由题设知，任一子列 $\{f_{k_i}(x)\}$ 中必有子列 $\{f_{k_{i_j}}(x)\}$ 在 $E$ 上几乎处处收敛于 $f(x)$。于是
\[
f_{k_{i_j}}^p(x)\to f^p(x)\qquad \aee\ x\in E.
\]
从而该子列也依测度收敛于 $f^p(x)$。根据 (i) 的子列判别准则，$f_k^p(x)$ 在 $E$ 上依测度收敛于 $f^p(x)$。

\section{第四章：Lebesgue 积分}

\subsection{非负可测函数积分}

\subsubsection*{题38（《实变函数论》p137 思考题3）}
\paragraph{题目}设 $\{f_k(x)\}$ 是 $E$ 上的非负可积函数列。若
\[
\lim_{k\to\infty}\int_E f_k(x)\,dx=0,
\]
证明
\[
\lim_{k\to\infty}\int_E\bigl(1-e^{-f_k(x)}\bigr)\,dx=0.
\]

\paragraph{答案}对 $t\ge0$，有基本不等式
\[
0\le 1-e^{-t}\le t.
\]
因此
\[
0\le\int_E\bigl(1-e^{-f_k(x)}\bigr)\,dx
\le\int_E f_k(x)\,dx.
\]
由题设右边趋于 $0$，故由夹逼定理，
\[
\int_E\bigl(1-e^{-f_k(x)}\bigr)\,dx\to0.
\]

\subsection{一般可积函数}

\subsubsection*{题39（《实变函数论》p154 思考题8）}
\paragraph{题目}设 $f\in L(\R)$。若对 $\R$ 上任意有界可测函数 $\varphi(x)$ 都有
\[
\int_{\R}f(x)\varphi(x)\,dx=0,
\]
证明
\[
f(x)=0,
\qquad \aee\ x\in\R.
\]

\paragraph{答案}取
\[
\varphi(x)=\operatorname{sgn} f(x)=
\begin{cases}
1,&f(x)>0,\\
0,&f(x)=0,\\
-1,&f(x)<0.
\end{cases}
\]
这是有界可测函数。由题设，
\[
0=\int_{\R}f(x)\operatorname{sgn} f(x)\,dx
=\int_{\R}|f(x)|\,dx.
\]
故 $|f|=0$ a.e.，即 $f=0$ a.e.。

\subsection{习题4中的积分题}

\subsubsection*{题40（《实变函数论》p190 习题5）}
\paragraph{题目}设 $f_k(x)$，$k=1,2,\ldots$，是 $\R^n$ 上非负可积函数列。若对任一可测集 $E\subset\R^n$，都有
\[
\int_E f_k(x)\,dx\le\int_E f_{k+1}(x)\,dx,
\]
证明
\[
\lim_{k\to\infty}\int_E f_k(x)\,dx
=
\int_E \lim_{k\to\infty}f_k(x)\,dx.
\]

\paragraph{答案}由题设对任意可测集 $E$ 成立，特别对集合
\[
E_k=\{x:f_k(x)>f_{k+1}(x)
\}
\]
成立。若 $\meas(E_k)>0$，则在 $E_k$ 上 $f_k-f_{k+1}>0$，从而
\[
\int_{E_k}f_k(x)\,dx>\int_{E_k}f_{k+1}(x)\,dx,
\]
矛盾。因此
\[
f_k(x)\le f_{k+1}(x),
\qquad \aee\ x.
\]
于是 $\{f_k\}$ 是非负渐升可测函数列。由 Beppo Levi 单调收敛定理，
\[
\lim_{k\to\infty}\int_E f_k(x)\,dx
=
\int_E \lim_{k\to\infty}f_k(x)\,dx.
\]

\subsubsection*{题41（《实变函数论》p190 习题7）}
\paragraph{题目}假设有定义在 $\R^n$ 上的函数 $f(x)$。如果对任意 $\eps>0$，存在 $g,h\in L(\R^n)$，满足
\[
g(x)\le f(x)\le h(x),
\qquad x\in\R^n,
\]
并且
\[
\int_{\R^n}(h(x)-g(x))\,dx<\eps,
\]
证明 $f\in L(\R^n)$。

\paragraph{答案}对每个 $m\in\N$，取 $g_m,h_m\in L(\R^n)$，使得
\[
g_m\le f\le h_m,
\qquad
\int_{\R^n}(h_m-g_m)<\frac1m.
\]
令
\[
u_m=\max\{g_1,
\ldots,g_m\},
\qquad
v_m=\min\{h_1,
\ldots,h_m\}.
\]
则 $u_m,v_m$ 可测，且
\[
u_m\le f\le v_m,
\qquad
0\le v_m-u_m\le h_m-g_m.
\]
于是
\[
\int_{\R^n}(v_m-u_m)\,dx\le\frac1m\to0.
\]
又 $u_m$ 递增，$v_m$ 递减，令
\[
u=\lim_{m\to\infty}u_m,
\qquad
v=\lim_{m\to\infty}v_m.
\]
则 $u\le f\le v$，且由 Fatou 引理或单调收敛定理可得
\[
\int_{\R^n}(v-u)\,dx=0.
\]
故 $u=v$ a.e.，于是 $f=u=v$ a.e.。又 $g_1\le u\le h_1$，而 $g_1,h_1\in L(\R^n)$，故 $u\in L(\R^n)$。因此 $f$ 与一个可积函数几乎处处相等，故 $f\in L(\R^n)$。

\subsubsection*{题42（《实变函数论》p190 习题10）}
\paragraph{题目}设 $f\in L(\R^n)$，$E\subset\R^n$ 是紧集。证明
\[
\lim_{|y|\to\infty}\int_{E+\{y\}}|f(x)|\,dx=0.
\]

\paragraph{答案}由于 $E$ 紧，存在 $R_0>0$，使得
\[
E\subset B(0,R_0).
\]
给定 $\eps>0$。因为 $f\in L(\R^n)$，存在 $R>0$，使得
\[
\int_{\{|x|>R\}}|f(x)|\,dx<\eps.
\]
当 $|y|>R+R_0$ 时，若 $x\in E+\{y\}$，则 $x=z+y$，其中 $z\in E$，于是
\[
|x|\ge |y|-|z|>R+R_0-R_0=R.
\]
故
\[
E+\{y\}\subset\{|x|>R\}.
\]
因此
\[
\int_{E+\{y\}}|f(x)|\,dx
\le
\int_{\{|x|>R\}}|f(x)|\,dx
<\eps.
\]
由 $\eps$ 任意，结论成立。

\newpage

\begin{center}
	{\Large \textbf{实变函数简明教程}}\\[0.3cm]
\end{center}

\section*{第一章 \quad 集合与点集\quad P38}

\subsubsection*{41题}
\paragraph{题目} (1) 证明：$\mathbb{R}^n$ 中只有 $\mathbb{R}^n$ 与空集 $\emptyset$ 是既开又闭的集合。\\
(2) 若 $E\subset\mathbb{R}^n$，$E\neq\emptyset$，$E\neq\mathbb{R}^n$，则边界 $\partial E\neq\emptyset$。

\paragraph{答案} (1) 显然 $\emptyset$ 与 $\mathbb{R}^n$ 都既开又闭。反过来，设 $E\subset\mathbb{R}^n$ 既开又闭，且 $E\neq\emptyset$。若 $E\neq\mathbb{R}^n$，取
\[
 x_0\in E,\qquad y_0\in E^c.
\]
考虑线段
\[
 \gamma(t)=(1-t)x_0+ty_0,\\qquad 0\le t\le 1.
\]
令
\[
 A=\{t\in[0,1]:\gamma(t)\in E\}.
\]
由于 $E$ 开，$A$ 是 $[0,1]$ 中的相对开集；由于 $E$ 闭，$A$ 又是 $[0,1]$ 中的相对闭集。并且 $0\in A$，$1\notin A$。这与区间 $[0,1]$ 的连通性矛盾。因此不存在非空真子集既开又闭，故 $\mathbb{R}^n$ 中只有 $\emptyset$ 和 $\mathbb{R}^n$ 既开又闭。

(2) 若 $\partial E=\emptyset$，则每个点要么是 $E$ 的内点，要么是 $E^c$ 的内点。因此 $E$ 与 $E^c$ 都是开集，即 $E$ 既开又闭。由 (1) 知 $E=\emptyset$ 或 $E=\mathbb{R}^n$，与题设矛盾。所以 $\partial E\neq\emptyset$。

\subsubsection*{49题}
\paragraph{题目} 设 $E\subset\mathbb{R}^n$。若 $E$ 的任意开球覆盖都存在有限子覆盖，证明 $E$ 是有界闭集。

\paragraph{答案} 先证有界性。任取 $x_0\in\mathbb{R}^n$，开球列
\[
 B(x_0,1),B(x_0,2),\ldots,B(x_0,k),\ldots
\]
覆盖 $E$。由题设可取有限个开球覆盖 $E$，于是存在 $N\in\mathbb{N}$，使得
\[
 E\subset B(x_0,N),
\]
故 $E$ 有界。

再证闭性。若 $E$ 不是闭集，则存在聚点 $y_0\in\overline E\setminus E$。由有界性，取 $N$ 使得 $E\subset B(x_0,N)$。对 $m=1,2,\ldots$，令
\[
 U_m=B(x_0,N)\setminus \overline{B\left(y_0,\frac1m\right)}.
\]
每个 $U_m$ 是开集，并且 $\{U_m\}_{m=1}^{\infty}$ 覆盖 $E$：事实上，对任意 $x\in E$，有 $x\neq y_0$，取 $m>1/|x-y_0|$ 即得 $x\in U_m$。

若存在有限子覆盖，令 $m_0$ 为其中最大指标，则由于 $U_m$ 随 $m$ 递增，有限子覆盖等价于 $E\subset U_{m_0}$。但 $y_0$ 是 $E$ 的聚点，故
\[
 E\cap B\left(y_0,\frac1{m_0}\right)\neq\emptyset,
\]
这与 $E\subset U_{m_0}$ 矛盾。因此 $E$ 必为闭集。综上，$E$ 是有界闭集。

\subsubsection*{50题}
\paragraph{题目} 设 $f$ 为 $\mathbb{R}$ 上连续函数，令
\[
 E=\{x\in\mathbb{R}:\exists\,t>x,\ f(t)>f(x)\}.
\]
证明：$E$ 是开集；且对 $E$ 的任意有限构成区间 $(a,b)$，有 $f(a)=f(b)$。

\paragraph{答案} 先证 $E$ 是开集。任取 $x\in E$，则存在 $t_0>x$，使得 $f(t_0)>f(x)$。由连续性，存在 $\delta>0$，并可取 $\delta<t_0-x$，使得当 $|y-x|<\delta$ 时，
\[
 f(y)<f(t_0).
\]
同时 $y<t_0$，故 $y\in E$。因此 $(x-\delta,x+\delta)\subset E$，所以 $E$ 为开集。

设 $(a,b)$ 是 $E$ 的一个有限构成区间。由于 $E$ 是开集，端点 $a,b\notin E$。由 $a\notin E$ 得
\[
 \forall t>a,\quad f(t)\le f(a),
\]
特别地 $f(b)\le f(a)$。又由 $b\notin E$ 得
\[
 \forall t>b,\quad f(t)\le f(b).
\]
下面证 $f(a)\le f(b)$。任取 $x\in(a,b)$。若 $f(x)>f(b)$，则由 $x\in E$，存在 $t>x$，使得 $f(t)>f(x)>f(b)$。此时 $t$ 不可能满足 $t\ge b$，因为 $t>b$ 时有 $f(t)\le f(b)$，而 $t=b$ 时也有 $f(t)=f(b)$。故必有 $t\in(x,b)$。

于是若某点 $x\in(a,b)$ 满足 $f(x)>f(b)$，则在右侧区间 $(x,b)$ 中还能找到函数值更大的点。取闭区间 $[x,b]$ 上 $f$ 的最大值点 $s$，若最大值大于 $f(b)$，则 $s\in[x,b)$，又 $s\in E$，故存在 $u>s$ 使 $f(u)>f(s)$，这与 $s$ 是 $[x,b]$ 上最大值点矛盾。因此对一切 $x\in(a,b)$ 都有
\[
 f(x)\le f(b).
\]
令 $x\to a+$，由连续性得 $f(a)\le f(b)$。结合 $f(b)\le f(a)$，即得
\[
 f(a)=f(b).
\]

\section*{第二章 \quad Lebesgue 测度\quad P55--P56}

\subsubsection*{3题}
\paragraph{题目} 设 $E\subset[0,1]$ 为可测集。若 $mE=1$，证明 $\overline E=[0,1]$；若 $mE=0$，证明 $E^\circ=\emptyset$。

\paragraph{答案} 若 $mE=1$。由于 $E\subset[0,1]$，故 $\overline E\subset[0,1]$。反设存在 $x_0\in[0,1]\setminus\overline E$，则存在 $\delta>0$，使得
\[
 (x_0-\delta,x_0+\delta)\cap[0,1]
\]
中含有一个非空开区间，且与 $E$ 不相交。于是 $[0,1]\setminus E$ 含有正测度子集，从而 $mE<1$，矛盾。因此 $[0,1]\subset\overline E$，即 $\overline E=[0,1]$。

若 $mE=0$。反设 $E^\circ\neq\emptyset$，则存在非空开区间 $I\subset E$。于是
\[
 0<mI\le mE=0,
\]
矛盾。故 $E^\circ=\emptyset$。

\subsubsection*{6题}
\paragraph{题目} 设 $E_1\subset E_2\subset\mathbb{R}^n$，$E_1$ 可测且 $mE_1=m^*E_2$，证明 $E_2$ 可测。

\paragraph{答案} 因为 $E_1$ 可测，故由可测集对任意集合的分割性质，
\[
 m^*E_2=m^*(E_2\cap E_1)+m^*(E_2\cap E_1^c).
\]
又 $E_1\subset E_2$，所以
\[
 E_2\cap E_1=E_1,\qquad E_2\cap E_1^c=E_2\setminus E_1.
\]
因此
\[
 m^*E_2=mE_1+m^*(E_2\setminus E_1).
\]
结合 $m^*E_2=mE_1$，得
\[
 m^*(E_2\setminus E_1)=0.
\]
零测集可测，故 $E_2\setminus E_1$ 可测。于是
\[
 E_2=E_1\cup(E_2\setminus E_1)
\]
为可测集之并，所以 $E_2$ 可测。

\subsubsection*{7题}
\paragraph{题目} 证明：有界集 $E$ 可测的充分必要条件是，对任意开集 $G$，都有
\[
 mG=m^*(G\cap E)+m^*(G\setminus E).
\]

\paragraph{答案} 必要性：若 $E$ 可测，则由 Carath\'{e}odory 可测条件，对任意集合 $A$ 都有
\[
 m^*A=m^*(A\cap E)+m^*(A\setminus E).
\]
特别取开集 $G$，且开集可测，$m^*G=mG$，便得
\[
 mG=m^*(G\cap E)+m^*(G\setminus E).
\]

充分性：任取 $A\subset\mathbb{R}^n$。若 $m^*A=+\infty$，则 Carath\'{e}odory 条件中的反向不等式自动成立。若 $m^*A<+\infty$，任给 $\eps>0$，按外测度定义可取开集 $G\supset A$，使得
\[
 mG<m^*A+\eps.
\]
由题设，
\[
 mG=m^*(G\cap E)+m^*(G\setminus E).
\]
又 $A\cap E\subset G\cap E$，$A\setminus E\subset G\setminus E$，故
\[
 m^*(A\cap E)+m^*(A\setminus E)
\le m^*(G\cap E)+m^*(G\setminus E)=mG<m^*A+\eps.
\]
令 $\eps\to0$，得
\[
 m^*(A\cap E)+m^*(A\setminus E)
\le m^*A.
\]
另一方面，外测度的次可加性总给出
\[
 m^*A\le m^*(A\cap E)+m^*(A\setminus E).
\]
故二者相等。由 Carath\'{e}odory 条件，$E$ 可测。

\subsubsection*{10题}
\paragraph{题目} 证明下列两个条件都是点集 $E$ 可测的充分必要条件：
\begin{enumerate}[label=(\arabic*)]
\item 对于任给的 $\eps>0$，存在开集 $G$ 和闭集 $F$，使得 $F\subset E\subset G$，且 $m(G\setminus F)<\eps$；
\item 对于任给的 $\eps>0$，存在开集 $G_1,G_2$，使得 $G_1\supset E$，$G_2\supset E^c$，且 $m(G_1\cap G_2)<\eps$。
\end{enumerate}

\paragraph{答案} 先证 (1)。若 $E$ 可测，则由可测集的开集逼近与闭集逼近性质，存在开集 $G\supset E$ 使
\[
 m(G\setminus E)<\frac{\eps}{2},
\]
又存在闭集 $F\subset E$ 使
\[
 m(E\setminus F)<\frac{\eps}{2}.
\]
于是
\[
 m(G\setminus F)\le m(G\setminus E)+m(E\setminus F)<\eps.
\]

反过来，若 (1) 成立，则对每个 $n\in\mathbb{N}$，可取闭集 $F_n$ 与开集 $G_n$，使
\[
 F_n\subset E\subset G_n,\qquad m(G_n\setminus F_n)<\frac1n.
\]
令
\[
 F=\bigcup_{n=1}^{\infty}F_n.
\]
则 $F$ 是 $F_\sigma$ 集，因而可测，且 $F\subset E$。并且对任意 $n$，
\[
 E\setminus F\subset E\setminus F_n\subset G_n\setminus F_n,
\]
所以
\[
 m^*(E\setminus F)\le \frac1n\qquad(n=1,2,\ldots).
\]
故 $m^*(E\setminus F)=0$，于是 $E=F\cup(E\setminus F)$ 可测。

再证 (2) 与 (1) 等价。若 (1) 成立，取
\[
 G_1=G,
\qquad
 G_2=F^c.
\]
则 $G_1\supset E$，$G_2\supset E^c$，且
\[
 G_1\cap G_2=G\cap F^c=G\setminus F,
\]
故 $m(G_1\cap G_2)<\eps$。

若 (2) 成立，取
\[
 G=G_1,
\qquad
 F=G_2^c.
\]
则 $F$ 为闭集，且由 $G_2\supset E^c$ 得 $F=G_2^c\subset E$；又 $G=G_1\supset E$。同时
\[
 G\setminus F=G_1\cap G_2,
\]
故 $m(G\setminus F)<\eps$。于是 (2) 推出 (1)。综上，(1)、(2) 都与 $E$ 可测等价。

\subsubsection*{11题}
\paragraph{题目} 设 $A,B\subset\mathbb{R}^n$ 都是可测集。证明在 $mA,mB<+\infty$ 的条件下，
\[
 m(A\cup B)=mA+mB
\]
的充分必要条件是 $A\cap B$ 是零测集。

\paragraph{答案} 先说明：若不加有限测度条件，原命题的必要性不成立。例如取 $A=B=\mathbb{R}^n$，则
\[
 m(A\cup B)=+\infty=mA+mB,
\]
但 $m(A\cap B)=+\infty$，并非零测集。因此这里按有限测度情形给出正确证明。

由于 $A,B$ 可测且测度有限，有
\[
 A\cup B=(A\setminus B)\cup B,
\]
且右端为不交并，故
\[
 m(A\cup B)=m(A\setminus B)+mB.
\]
又
\[
 A=(A\setminus B)\cup(A\cap B)
\]
为不交并，所以
\[
 mA=m(A\setminus B)+m(A\cap B).
\]
因此
\[
 mA+mB=m(A\cup B)+m(A\cap B).
\]
于是
\[
 m(A\cup B)=mA+mB
\]
当且仅当
\[
 m(A\cap B)=0.
\]

\subsubsection*{12题}
\paragraph{题目} 设 $f$ 是定义在 $\mathbb{R}$ 上的实值连续函数，证明它的图像
\[
 \Gamma_f=\{(x,y)\in\mathbb{R}^2:y=f(x),\ x\in\mathbb{R}\}
\]
是 $\mathbb{R}^2$ 中的零测集。

\paragraph{答案} 对每个 $N\in\mathbb{N}$，记
\[
 \Gamma_N=\{(x,f(x)):-N\le x\le N\}.
\]
由于
\[
 \Gamma_f=\bigcup_{N=1}^{\infty}\Gamma_N,
\]
只需证明每个 $\Gamma_N$ 是二维零测集。

固定 $N$。函数 $f$ 在紧区间 $[-N,N]$ 上一致连续。任给 $\eps>0$，取 $\eta>0$，使得
\[
 4N\eta<\eps.
\]
由一致连续性，存在 $\delta>0$，使得当 $|x-x'|<\delta$ 时，
\[
 |f(x)-f(x')|<\eta.
\]
将 $[-N,N]$ 分成有限个小区间 $I_j=[x_{j-1},x_j]$，并使每个小区间长度小于 $\delta$。对每个 $I_j$，取一点 $\xi_j\in I_j$，则当 $x\in I_j$ 时，
\[
 |f(x)-f(\xi_j)|<\eta.
\]
于是 $\Gamma_N$ 被有限个矩形
\[
 I_j\times(f(\xi_j)-\eta,\ f(\xi_j)+\eta)
\]
覆盖。这些矩形面积总和不超过
\[
 \sum_j |I_j|\cdot 2\eta=2\eta\cdot 2N=4N\eta<\eps.
\]
由 $\eps$ 任意，$m_2^*(\Gamma_N)=0$。故每个 $\Gamma_N$ 为二维零测集，从而 $\Gamma_f$ 是可数个零测集之并，仍为零测集。

\section*{第三章 \quad 可测函数\quad P78--P79}

\subsubsection*{3题}
\paragraph{题目} 若每个 $f_k(x)$ 在可测集 $E\subset\mathbb{R}^n$ 上几乎处处连续，即间断点构成零测集，且极限
\[
 f(x)=\lim_{k\to\infty}f_k(x)
\]
在 $E$ 上几乎处处有意义，证明 $f$ 在 $E$ 上可测。

\paragraph{答案} 对每个 $k$，设 $N_k$ 是 $f_k$ 在 $E$ 中的间断点集，则 $mN_k=0$。在 $E\setminus N_k$ 上，$f_k$ 连续，因而可测；在零测集 $N_k$ 上任意修改函数值不影响 Lebesgue 可测性，所以 $f_k$ 在 $E$ 上可测。

又因为 $f(x)=\lim_{k\to\infty}f_k(x)$ 在 $E$ 上几乎处处有意义，存在零测集 $N_0\subset E$，使得在 $E\setminus N_0$ 上处处有
\[
 f(x)=\lim_{k\to\infty}f_k(x).
\]
可测函数列的逐点极限仍可测，故 $f$ 在 $E\setminus N_0$ 上可测。再在零测集 $N_0$ 上任意定义或修改 $f$，仍不影响可测性，所以 $f$ 在 $E$ 上可测。

\subsubsection*{4题}
\paragraph{题目} 若 $f(x)$ 是集 $E$ 上的可测函数，证明对于任意实数 $a$，$E(f=a)$ 是可测集。反过来，若对任意 $a\in\mathbb{R}$，$E(f=a)$ 可测，能否断言 $f$ 可测？若还知道 $f(x)$ 只取可数个不同的值，$f$ 可测吗？

\paragraph{答案} 若 $f$ 可测，则
\[
 E(f=a)=E(f\ge a)\setminus E(f>a),
\]
右端为可测集之差，所以 $E(f=a)$ 可测。

反过来，不能断言 $f$ 可测。取不可测集 $A\subset[-\frac12,\frac12]$，令 $E=[-\frac12,\frac12]$，定义
\[
 f(x)=
 \begin{cases}
 x, & x\in A,\\
 x-2, & x\in E\setminus A.
 \end{cases}
\]
则对任意 $a\in\mathbb{R}$，集合 $E(f=a)$ 至多为单点集，故可测。但
\[
 E(f>-1)=A,
\]
不可测，所以 $f$ 不可测。

若 $f$ 只取可数个不同的值，设这些值为 $a_1,a_2,\ldots$。对任意 $c\in\mathbb{R}$，有
\[
 E(f>c)=\bigcup_{a_j>c}E(f=a_j),
\]
这是可数个可测集之并，故可测。因此 $f$ 可测。

\subsubsection*{7题}
\paragraph{题目} 设函数 $f$ 是在有界集 $E\subset\mathbb{R}^n$ 上几乎处处有限的可测函数。
\begin{enumerate}[label=(\arabic*)]
\item 证明：对于任意正数 $\delta$，存在 $E$ 的可测子集 $E_0$ 与正数 $M$，使得 $m(E\setminus E_0)<\delta$ 且 $|f(x)|<M$，$x\in E_0$；
\item 是否一定存在正数 $M$ 和 $E$ 的可测子集 $E_0$，使得 $m(E\setminus E_0)=0$ 且 $|f(x)|<M$，$x\in E_0$？
\end{enumerate}

\paragraph{答案} (1) 因为 $E$ 有界，所以 $mE<+\infty$。又 $f$ 几乎处处有限，故
\[
 E_N=E(|f|<N)\qquad(N=1,2,\ldots)
\]
是递增可测集列，且
\[
 \bigcup_{N=1}^{\infty}E_N=E\setminus E(|f|=+\infty),
\]
其中 $mE(|f|=+\infty)=0$。由测度连续性，
\[
 m(E\setminus E_N)\to0.
\]
故对给定 $\delta>0$，可取 $N$ 充分大，使
\[
 m(E\setminus E_N)<\delta.
\]
令 $E_0=E_N$，$M=N$，即得结论。

(2) 不一定。例如取 $E=(0,1)$，$f(x)=1/x$。对任意 $M>0$，集合
\[
 \{x\in(0,1):|f(x)|\ge M\}
\]
都含有 $(0,1/M]$ 的一部分，从而有正测度。因此不可能删去一个零测集后使 $f$ 有界。

\subsubsection*{9题（1）（2）}
\paragraph{题目} (1) 若 $f$ 是集 $E\subset\mathbb{R}^n$ 上的实值函数，证明：对于任意开集 $G\subset\mathbb{R}$ 或闭集 $F\subset\mathbb{R}$，点集
\[
 \{x\in E:f(x)\in G\},\qquad \{x\in E:f(x)\in F\}
\]
都可测，是 $f$ 在 $E$ 上可测的充要条件。\\
(2) 设 $f$ 是可测集 $E\subset\mathbb{R}^n$ 上的广义实值函数。证明 $f$ 为可测函数的充要条件是：
\begin{enumerate}[label=(\roman*)]
\item $E(f=-\infty)$ 与 $E(f=+\infty)$ 都是可测集；
\item 对于 $\mathbb{R}$ 中的任意 Borel 集 $B$，集合 $\{x\in E:f(x)\in B\}$ 总是可测集。
\end{enumerate}

\paragraph{答案} (1) 必要性：若 $f$ 可测，则对任意开区间 $(a,b)$，
\[
 E(a<f<b)=E(f>a)\cap E(f<b)
\]
可测。任意开集 $G\subset\mathbb{R}$ 可写成至多可数个互不相交开区间的并，故
\[
 \{x\in E:f(x)\in G\}
\]
可测。若 $F$ 是闭集，则
\[
 \{x\in E:f(x)\in F\}=E\setminus\{x\in E:f(x)\in F^c\},
\]
也可测。

充分性：若对任意开集原像都可测，则特别取 $G=(a,+\infty)$，得到
\[
 E(f>a)=\{x\in E:f(x)\in(a,+\infty)
\}
\]
可测。故 $f$ 可测。

(2) 必要性：若 $f$ 是广义实值可测函数，则
\[
 E(f=+\infty)=\bigcap_{N=1}^{\infty}E(f>N),
\qquad
E(f=-\infty)=E\setminus\bigcup_{N=1}^{\infty}E(f>-N),
\]
均可测。又在
\[
 E_0=E\setminus(E(f=+\infty)\cup E(f=-\infty))
\]
上，$f$ 是实值可测函数。由 (1) 以及 Borel 集由开集生成，可知对任意 Borel 集 $B\subset\mathbb{R}$，
\[
 \{x\in E:f(x)\in B\}=\{x\in E_0:f(x)\in B\}
\]
可测。

充分性：若 (i)(ii) 成立，则对任意 $a\in\mathbb{R}$，由于 $(a,+\infty)$ 是 Borel 集，
\[
 E(f>a)=\{x\in E:a<f(x)<+\infty\}\cup E(f=+\infty)
\]
可测。故 $f$ 是广义实值可测函数。

\subsubsection*{15题}
\paragraph{题目} 设 $f_k(x)$ 在 $E$ 上依测度收敛于 $f(x)$，$g_k(x)$ 在 $E$ 上依测度收敛于 $g(x)$，且 $mE<\infty$。证明
\[
 f_k(x)g_k(x)\to f(x)g(x)
\]
在 $E$ 上依测度成立。

\paragraph{答案} 先用一个简单事实：若 $u_k\to0$ 依测度，$v_k\to v$ 依测度，且 $mE<\infty$，则 $u_kv_k\to0$ 依测度。证明如下：任给 $\eta>0$ 与 $\eps>0$。由于 $v$ 几乎处处有限且 $mE<\infty$，可取 $M>0$ 使
\[
 mE(|v|>M)<\eps.
\]
又由 $v_k\to v$ 依测度，
\[
 mE(|v_k-v|>1)\to0.
\]
因此当 $k$ 充分大时，
\[
 mE(|v_k|>M+1)
\le mE(|v|>M)+mE(|v_k-v|>1)<2\eps.
\]
于是
\[
 E(|u_kv_k|>\eta)
\subset E(|v_k|>M+1)\cup E\left(|u_k|>\frac{\eta}{M+1}\right).
\]
令 $k\to\infty$ 得
\[
 \limsup_{k\to\infty}mE(|u_kv_k|>\eta)\\le 2\eps.
\]
再令 $\eps\to0$，即得 $u_kv_k\to0$ 依测度。

现在分解
\[
 f_kg_k-fg=(f_k-f)g_k+f(g_k-g).
\]
由上述事实，取 $u_k=f_k-f$、$v_k=g_k$，得
\[
 (f_k-f)g_k\to0
\]
依测度；再取 $u_k=g_k-g$、$v_k\equiv f$，得
\[
 f(g_k-g)\to0
\]
依测度。两个依测度收敛到 $0$ 的函数列之和仍依测度收敛到 $0$，故
\[
 f_kg_k-fg\to0
\]
依测度，即 $f_kg_k\to fg$ 依测度。

\section*{第四章 \quad Lebesgue 积分 \quad P114--P116}

\subsubsection*{2题}
\paragraph{题目} 若 $mE<\infty$，$f$ 是在 $E$ 上几乎处处有限的非负可测函数，证明
\[
 \int_E f(x)\,dx<+\infty
\]
的充分必要条件是
\[
 \sum_{k=1}^{\infty} k\,mE(k\le f<k+1)<+\infty.
\]

\paragraph{答案} 记
\[
 A_k=E(k\le f<k+1),\qquad k=1,2,\ldots.
\]
由于 $f$ 几乎处处有限，$E(f=+\infty)$ 为零测集。对 $x\in A_k$，有
\[
 k\le f(x)<k+1.
\]
因此
\[
 \sum_{k=1}^{\infty} k\,mA_k
\le \int_E f(x)\,dx.
\]
所以若积分有限，则级数收敛。

反过来，若
\[
 \sum_{k=1}^{\infty} k\,mA_k<+\infty,
\]
则由 $mE<\infty$，
\[
 \sum_{k=1}^{\infty}mA_k\le mE<+\infty.
\]
于是
\[
 \sum_{k=1}^{\infty}(k+1)mA_k<+\infty.
\]
又在 $E(0\le f<1)$ 上有 $f\le1$，在 $A_k$ 上有 $f<k+1$，故
\[
 \int_E f(x)\,dx
\le mE(0\le f<1)+\sum_{k=1}^{\infty}(k+1)mA_k<+\infty.
\]
证毕。

\subsubsection*{3题}
\paragraph{题目} 若 $f$ 是测度有限的点集 $E$ 上的非负可测函数，证明
\[
 \int_E f(x)\,dx<+\infty
\]
的充分必要条件是
\[
 \sum_{k=1}^{\infty}mE(f\ge k)<+\infty.
\]

\paragraph{答案} 对每个 $x\in E$，有点态估计
\[
 \sum_{k=1}^{\infty}\chi_{E(f\ge k)}(x)\le f(x).
\]
事实上，左端等于不超过 $f(x)$ 的正整数个数，即 $\lfloor f(x)\rfloor$，自然不超过 $f(x)$。由单调收敛定理，若 $\int_E f<+\infty$，则
\[
 \sum_{k=1}^{\infty}mE(f\ge k)
=\int_E\sum_{k=1}^{\infty}\chi_{E(f\ge k)}(x)\,dx
\le \int_E f(x)\,dx<+\infty.
\]

反过来，若
\[
 \sum_{k=1}^{\infty}mE(f\ge k)<+\infty,
\]
则对每个 $x\in E$，有
\[
 f(x)\le 1+\sum_{k=1}^{\infty}\chi_{E(f\ge k)}(x),
\]
这里在 $0\le f(x)<1$ 时右端为 $1$，在 $f(x)\ge1$ 时右端为 $1+\lfloor f(x)\rfloor$，足以控制 $f(x)$；若 $f(x)=+\infty$，右端也为 $+\infty$。因此
\[
 \int_E f(x)\,dx
\le mE+\sum_{k=1}^{\infty}mE(f\ge k)<+\infty.
\]
证毕。


\subsubsection*{16题}
\paragraph{题目} 设 $f$ 是实直线 $\mathbb{R}$ 上的可积函数，$f(0)=0$，导数 $f'(0)$ 存在且有限。试证
\[
 \frac{f(x)}{x}\in L(\mathbb{R}).
\]
这里 $x=0$ 处的函数值可任意规定，不影响 Lebesgue 可积性。

\paragraph{答案} 记
\[
 h(x)=\begin{cases}
 \dfrac{f(x)}{x},&x\ne0,\\[4pt]
 f'(0),&x=0.
 \end{cases}
\]
只需证明 $h\in L(\mathbb{R})$。由于 $f'(0)$ 存在且有限，由导数定义，存在 $\delta>0$，使得当 $0<|x|<\delta$ 时，
\[
 \left|\frac{f(x)-f(0)}{x-0}-f'(0)\right|<1.
\]
又 $f(0)=0$，所以当 $0<|x|<\delta$ 时，
\[
 \left|\frac{f(x)}{x}\right|\le |f'(0)|+1.
\]
于是
\[
 \int_{|x|<\delta}\left|\frac{f(x)}{x}\right|\,dx
 \le 2\delta\bigl(|f'(0)|+1\bigr)<+\infty.
\]
另一方面，在 $|x|\ge\delta$ 上有 $1/|x|\le1/\delta$，故
\[
 \int_{|x|\ge\delta}\left|\frac{f(x)}{x}\right|\,dx
 \le \frac1\delta\int_{\mathbb{R}}|f(x)|\,dx<+\infty,
\]
因为 $f\in L(\mathbb{R})$。综上，
\[
 \int_{\mathbb{R}}\left|\frac{f(x)}{x}\right|\,dx
 =\int_{|x|<\delta}\left|\frac{f(x)}{x}\right|\,dx
 +\int_{|x|\ge\delta}\left|\frac{f(x)}{x}\right|\,dx<+\infty.
\]
因此 $\dfrac{f(x)}{x}\in L(\mathbb{R})$。


\subsubsection*{26题}
\paragraph{题目} 设 $mE<\infty$，$\{f_n\}$ 是集 $E$ 上的可测函数列，试证当 $n\to\infty$ 时 $f_n$ 依测度收敛于 $0$ 的充分必要条件是
\[
 \lim_{n\to\infty}\int_E\frac{|f_n(x)|}{1+|f_n(x)|}\,dx=0.
\]

\paragraph{答案} 令
\[
 \Phi(t)=\frac{t}{1+t},\qquad t\ge0.
\]
则 $0\le\Phi(t)<1$，且 $\Phi$ 在 $[0,+\infty)$ 上单调递增。

先证必要性。设 $f_n\to0$ 依测度于 $E$。若 $mE=0$，结论显然；下设 $mE>0$。任取 $\varepsilon>0$，取 $0<\eta<1$，使
\[
 \eta mE<\frac{\varepsilon}{2}.
\]
由依测度收敛，
\[
 mE(|f_n|>\eta)\to0.
\]
故当 $n$ 充分大时，
\[
 mE(|f_n|>\eta)<\frac{\varepsilon}{2}.
\]
把积分按集合 $E(|f_n|>\eta)$ 与 $E(|f_n|\le\eta)$ 分开，得
\[
\begin{aligned}
 \int_E\frac{|f_n(x)|}{1+|f_n(x)|}\,dx
 &=\int_{E(|f_n|>\eta)}\frac{|f_n(x)|}{1+|f_n(x)|}\,dx
   +\int_{E(|f_n|\le\eta)}\frac{|f_n(x)|}{1+|f_n(x)|}\,dx \\
 &\le mE(|f_n|>\eta)+\eta mE
 <\frac{\varepsilon}{2}+\frac{\varepsilon}{2}=\varepsilon.
\end{aligned}
\]
所以
\[
 \int_E\frac{|f_n(x)|}{1+|f_n(x)|}\,dx\to0.
\]

再证充分性。设
\[
 \int_E\frac{|f_n(x)|}{1+|f_n(x)|}\,dx\to0.
\]
任取 $\varepsilon>0$。在集合 $E(|f_n|\ge\varepsilon)$ 上，由 $\Phi$ 的单调性，
\[
 \frac{|f_n(x)|}{1+|f_n(x)|}\ge \frac{\varepsilon}{1+\varepsilon}.
\]
因此
\[
 \int_E\frac{|f_n(x)|}{1+|f_n(x)|}\,dx
 \ge \int_{E(|f_n|\ge\varepsilon)}\frac{|f_n(x)|}{1+|f_n(x)|}\,dx
 \ge \frac{\varepsilon}{1+\varepsilon}mE(|f_n|\ge\varepsilon).
\]
于是
\[
 mE(|f_n|\ge\varepsilon)
 \le \frac{1+\varepsilon}{\varepsilon}
 \int_E\frac{|f_n(x)|}{1+|f_n(x)|}\,dx\to0.
\]
这正是 $f_n\to0$ 依测度于 $E$。故充分必要性得证。

\end{document}
